% Part: model-theory
% Chapter: models-of-arithmetic
% Section: standard-models

\documentclass[../../../include/open-logic-section]{subfiles}

\begin{document}

\olfileid{mod}{mar}{stm}
\olsection{حساب دے معیاری نمونے}

حساب دی زبان~$\Lang{L_A}$ دا مقصد ظاہر اے کہ اعداد،
خاص طور اُتے قدرتی اعداد، بارے گل کرنا اے۔ ایس لئی ”اوہ“
معیاری نمونہ~$\Struct{N}$ خاص اے: ایہ اوہ نمونہ اے
جیہدے بارے اسیں گل کرنا چاہندے آں۔ پر منطق وچ ساڈی
دلچسپی اکثر ساختی خاصیتاں وچ ہوندی اے، تے ہر دو
ہم شکل !!{structure}s ایہ خاصیتاں سانجھیاں رکھدیاں نیں۔
ایس لئی اسیں تعریف ذرا کھلی رکھ سکدے آں تے ہر اوہ
!!{structure} ”معیاری“ من سکدے آں جیہڑی~$\Struct{N}$
نال ہم شکل ہووے۔

\begin{defn}
$\Lang{L_A}$ لئی !!{structure} \emph{معیاری} اے
جے اوہ~$\Struct{N}$ نال ہم شکل ہووے۔
\end{defn}

\begin{prop}
\ollabel{prop:standard-domain} جے !!a{structure}~$\Struct{M}$
معیاری اے، تاں اوہدا دائرہ معیاری عددی عبارتاں دیاں
قدراں دا سیٹ اے، یعنی
\[
\Domain{M} = \Setabs{\Value{\num{n}}{M}}{n \in \Nat}
\]
\end{prop}

\begin{proof}
واضح اے کہ ہر $\Value{\num{n}}{M} \in \Domain{M}$۔
بس ایہ وکھانا اے کہ ہر $x \in \Domain{M}$
$\Value{\num{n}}{M}$ دے برابر اے، کسے~$n$ لئی۔
کیوں جے $\Struct{M}$ معیاری اے، ایہ~$\Struct{N}$
نال ہم شکل اے۔ فرض کرو $g\colon \Nat \to \Domain{M}$
اک ہم شکلی اے۔ فیر $g(n) = g(\Value{\num{n}}{N}) =
\Value{\num{n}}{M}$۔ پر ہر $x \in \Domain{M}$
لئی کوئی~$n \in \Nat$ اے جیہدے لئی $g(n) = x$،
کیوں جے $g$ ہر ہدف قدر لیندا اے (!!{surjective})۔
\end{proof}

\begin{explain}
جے ساخت~$\Struct{M}$، $\Lang{L_A}$ لئی، معیاری اے،
تاں اوہدے !!{domain} دے سارے عنصراں نوں
معیاری عددی عبارتاں $\num{0}$، $\num{1}$،
$\num{2}$، \dots نال ناں دتا جا سکدا اے؛
یعنی ٹرماں $\Obj{0}$، $\Obj{0}'$،
$\Obj{0}''$، وغیرہ نال۔ بے شک ایس دا ایہ مطلب نہیں کہ
$\Domain{M}$ دے !!{element}s خود اعداد \emph{نیں}؛
صرف ایہ کہ اسیں اوہناں نوں اوہو طرح چن سکدے آں
جیویں~$\Domain{N}$ وچ اعداد نوں چن سکدے آں۔
\end{explain}

\begin{prob}
وکھاؤ کہ \olref[mod][mar][stm]{prop:standard-domain}
دا الٹا دعویٰ غلط اے؛ یعنی ایسی !!a{structure}~$\Struct{M}$
دی مثال دیو جیہدے لئی
$\Domain{M} = \Setabs{\Value{\num{n}}{M}}{n \in \Nat}$
ہووے، پر اوہ~$\Struct{N}$ نال ہم شکل نہ ہووے۔
\end{prob}

\begin{prop}
\ollabel{prop:thq-standard}
جے $\Sat{M}{\Th{Q}}$ ہووے تے
$\Domain{M} = \Setabs{\Value{\num{n}}{M}}{n
  \in \Nat}$، تاں $\Struct{M}$ معیاری اے۔
\end{prop}

\begin{proof}
اسیں وکھانا اے کہ $\Struct{M}$،
~$\Struct{N}$ نال ہم شکل اے۔ فنکشن
$g\colon \Nat \to \Domain{M}$ ویکھو جیہدی تعریف
$g(n) = \Value{\num{n}}{M}$ اے۔ مفروضے مطابق
$g$ ہر ہدف قدر لیندا اے (!!{surjective})۔ ایہ
اک بہ اک (!!{injective}) وی اے:
$\Th{Q} \Proves
\eq/[\num{n}][\num{m}]$ جدوں $n \neq m$
ہووے۔ ایس لئی، کیوں جے
$\Sat{M}{\Th{Q}}$، تاں
$\Sat{M}{\eq/[\num{n}][\num{m}]}$،
جدوں $n \neq m$۔ ایس لئی جے $n \neq m$، تاں
$\Value{\num{n}}{M} \neq
\Value{\num{m}}{M}$، یعنی $g(n) \neq g(m)$۔

ایہ وی ثابت کرنا اے کہ $g$ ہم شکلی اے۔
\begin{enumerate}
\item ساڈے کول $g(\Assign{\Obj{0}}{N}) = g(0)$ اے،
  کیوں جے $\Assign{\Obj{0}}{N} = 0$۔ $g$ دی تعریف توں
  $g(0) = \Value{\num{0}}{M}$۔ پر $\num{0}$
  بس $\Obj{0}$ اے، تے جیہڑا ٹرم !!a{constant} ہووے
  اوہدی قدر اوہ اے جیہڑی !!{structure} اوس
  !!{constant} نوں دندی اے، یعنی
  $\Value{\Obj{0}}{M} = \Assign{\Obj{0}}{M}$۔
  ایس لئی مطلوبہ برابری
  $g(\Assign{\Obj{0}}{N}) = \Assign{\Obj{0}}{M}$ مل گئی۔
\item $g(\Assign{\prime}{N}(n)) = g(n+1)$،
  کیوں جے $\prime$، $\Struct{N}$ وچ
  ~$\Nat$ اُتے جانشین فنکشن اے۔ فیر
  $g(n+1) = \Value{\num{n+1}}{M}$،
  $g$ دی تعریف توں۔
  پر $\num{n+1}$ اوہو ٹرم اے جو $\num{n}'$،
  ایس لئی $\Value{\num{n+1}}{M} =
  \Value{\num{n}'}{M}$۔ قدر والے فنکشن دی تعریف توں
  ایہ $= \Assign{\prime}{M}(\Value{\num{n}}{M})$ اے۔
  کیوں جے $\Value{\num{n}}{M} = g(n)$،
  اسیں $g(\Assign{\prime}{N}(n)) =
  \Assign{\prime}{M}(g(n))$ حاصل کردے آں۔
\item $g(\Assign{+}{N}(n,m)) = g(n+m)$،
  کیوں جے $+$، $\Struct{N}$ وچ
  ~$\Nat$ اُتے جمع دا فنکشن اے۔ فیر
  $g(n+m) = \Value{\num{n+m}}{M}$،
  $g$ دی تعریف توں۔
  پر $\Th{Q} \Proves
  \num{n+m} = (\num{n} + \num{m})$،
  ایس لئی $\Value{\num{n+m}}{M} =
  \Value{\num{n}+\num{m}}{M}$۔ قدر والے فنکشن دی
  تعریف توں ایہ $=
  \Assign{+}{M}(\Value{\num{n}}{M},\Value{\num{m}}{M})$
  اے۔ کیوں جے $\Value{\num{n}}{M} = g(n)$ تے
  $\Value{\num{m}}{M} = g(m)$، اسیں
  $g(\Assign{+}{N}(n, m)) = \Assign{+}{M}(g(n), g(m))$
  حاصل کردے آں۔
\item $g(\Assign{\times}{N}(n, m)) =
  \Assign{\times}{M}(g(n), g(m))$: مشق۔
\item $\tuple{n,m} \in \Assign{<}{N}$ اُدوں تے صرف اُدوں
  جدوں $n < m$۔ جے $n < m$، تاں
  $\Th{Q} \Proves \num{n} < \num{m}$،
  تے $\Sat{M}{\num{n} <
  \num{m}}$ وی۔ ایس لئی
  $\tuple{\Value{\num{n}}{M}, \Value{\num{m}}{M}} \in
  \Assign{<}{M}$، یعنی
  $\tuple{g(n), g(m)} \in \Assign{<}{M}$۔
  جے $n \not< m$، تاں
  $\Th{Q} \Proves \lnot \num{n} < \num{m}$،
  تے نتیجتاً $\Sat/{M}{\num{n} < \num{m}}$۔
  ایس لئی، پہلے وانگ،
  $\tuple{g(n), g(m)} \notin \Assign{<}{M}$۔
  دوہاں حالاں نوں ملا کے
  $\tuple{n,m} \in \Assign{<}{N}$
  اُدوں تے صرف اُدوں جدوں
  $\tuple{g(n), g(m)} \in
  \Assign{<}{M}$۔
\end{enumerate}
\end{proof}

\begin{explain}
فنکشن~$g$، $\Nat$ توں کسے وی ہور
!!{structure}~$\Struct{M}$ دے دائرے ول، $\Lang{L_A}$
لئی، نقشہ متعین کرن
دا سب توں سیدھا طریقہ اے، کیوں جے ہر ایسی
$\Struct{M}$ وچ $\num{0}$، $\num{1}$،
$\num{2}$ وغیرہ نال ناں ملن والے !!{element}s
ہوندے نیں۔ ایس لئی حیرانی نہیں کہ جے $\Struct{M}$،
$\num{n}$ دے بارے گھٹ توں گھٹ کجھ بنیادی بیان
~$\Struct{N}$ وانگ سچے بناؤندی اے، تے $g$ اک بہ اک
مطابقت وی اے، تاں $g$ ہم شکلی بن جاندا اے۔
ایہناں مفروضیاں تحت، جے $\Domain{M}$ وچ $\num{n}$
دے ناں والے عنصراں توں سوا کوئی !!{element}s نہیں،
تاں ایہی واحد ہم شکلی اے۔
\end{explain}

\begin{prop}
\ollabel{prop:thq-unique-iso} جے $\Struct{M}$ معیاری اے،
تاں \olref{prop:thq-standard} دے ثبوت والا $g$،
$\Struct{N}$ توں~$\Struct{M}$ ول واحد ہم شکلی اے۔
\end{prop}

\begin{proof}
فرض کرو $h\colon \Nat \to \Domain{M}$،
$\Struct{N}$ تے~$\Struct{M}$ وچ ہم شکلی اے۔
اسیں $g = h$ نوں~$n$ اُتے استقرا نال وکھاؤندے آں۔
جے $n = 0$، تاں
$g(0) = \Assign{\Obj{0}}{M}$، $g$ دی تعریف توں۔
پر کیوں جے $h$
ہم شکلی اے، $h(0) =
h(\Assign{\Obj{0}}{N}) =\Assign{\Obj{0}}{M}$،
ایس لئی $g(0) = h(0)$۔

ہُن $n+1$ دی حالت ویکھو۔ ساڈے کول
\begin{align*}
  g(n+1) & = \Value{\num{n+1}}{M} \text{ $g$ دی تعریف توں}\\
  & = \Value{\num{n}'}{M} \text{ کیوں جے $\num{n+1}\ident \num{n}'$}\\
  & = \Assign{\prime}{M}(\Value{\num{n}}{M})
    \text{ $\Value{t'}{M}$ دی تعریف توں}\\
  & = \Assign{\prime}{M}(g(n))  \text{ $g$ دی تعریف توں}\\
  & = \Assign{\prime}{M}(h(n)) \text{ استقرائی مفروضے توں}\\
  & = h(\Assign{\prime}{N}(n)) \text{ کیوں جے $h$ ہم شکلی اے}\\
  & = h(n+1)
\end{align*}
\end{proof}

\begin{explain}
ہر !!{denumerable} سیٹ~$M$ لئی $\Nat$ تے $M$
وچکار !!a{bijection} ہوندی اے، ایس لئی ہر ایسا سیٹ~$M$
کسی معیاری نمونے~$\Struct{M}$ دا !!{domain}
بن سکدا اے۔ بلکہ جدوں اک شے $z \in
M$ تے اک مناسب فنکشن $s$، بالترتیب
$\Assign{\Obj{0}}{M}$ تے
$\Assign{\prime}{M}$ دے طور اُتے چن لو،
تاں $+$، $\times$ تے $<$ دیاں تعبیراں
پہلے ای مقرر ہو جاندیاں نیں۔ صرف اوہ فنکشن
$s\colon M \to M \setminus \{z\}$،
معیاری نمونے وچ~$\Assign{\prime}{M}$ لئی مناسب نیں
جو اک بہ اک (!!{injective}) تے ہر ہدف قدر لین والے
(!!{surjective}) دوویں ہون۔ $s$ دے حاصل قدراں دے سیٹ وچ~$z$
نہیں آ سکدا؛ نہیں تاں
$\lforall[x][\eq/[\Obj 0][x']]$ غلط ہوندا۔
ایہ !!{sentence}، $\Struct{N}$ وچ سچ اے،
تے ایس لئی $\Struct{M}$ وچ وی سچ ہونا چاہیدا اے۔
فنکشن~$s$ دا اک بہ اک (!!{injective}) ہونا لازم اے،
کیوں جے جانشین فنکشن~$\Assign{\prime}{N}$،
$\Struct{N}$ وچ ایداں دا اے، تے
~$\Assign{\prime}{N}$ دا اک بہ اک (!!{injective}) ہونا
!!a{sentence} نال بیان کیتا جاندا اے جیہڑا
~$\Struct{N}$ وچ سچ اے۔ اوہدا ہر ہدف قدر لین والا
(!!{surjective}) ہونا وی لازم اے، نہیں تاں کوئی
$x \in M \setminus \{z\}$،
~$s$ دے حاصل قدراں دے سیٹ وچ نہ ہوندا؛ یعنی
!!{sentence} $\lforall[x][(\eq[x][\Obj 0] \lor
\lexists[y][\eq[y'][x]])]$،
~$\Struct{M}$ وچ غلط ہوندا، حالانکہ
~$\Struct{N}$ وچ سچ اے۔
\end{explain}


\end{document}
